\documentclass[11pt,a4paper]{article}
\usepackage[utf8]{inputenc}
\usepackage[T1]{fontenc}
\usepackage{amsmath,amssymb,amsthm,physics,geometry,booktabs,array,hyperref,mathtools}
\geometry{margin=2.5cm}
\hypersetup{colorlinks=true,linkcolor=blue,citecolor=blue,urlcolor=blue}
\theoremstyle{plain}
\newtheorem{theorem}{Theorem}[section]
\newtheorem{lemma}[theorem]{Lemma}
\newtheorem{corollary}[theorem]{Corollary}
\newtheorem{proposition}[theorem]{Proposition}
\theoremstyle{definition}
\newtheorem{definition}{Definition}[section]
\theoremstyle{remark}
\newtheorem{remark}{Remark}[section]
\newcommand{\lQ}{\ell_Q}
\newcommand{\Ms}{M_\odot}
\newcommand{\boldx}{\mathbf{x}}
\newcommand{\bhat}[1]{\hat{\mathbf{#1}}}
\newcommand{\sig}[2]{\sigma_{#1}^{(#2)}}
\newcommand{\rs}[1]{r_{\mathrm{s}}^{(#1)}}
\newcommand{\crossterm}[2]{\sigma_t^{(#1)}\sigma_t^{(#2)}}
\newcommand{\Sigt}{\Sigma_{\mathrm{tot}}}
\newcommand{\Fp}{F_{\mathrm{QGD}}^{\mathrm{dip}}}

\title{\textbf{Quantum Gravitational Dynamics}\\[4pt]
\Large Chapter 6: Exact Solutions and the Post-Newtonian Hierarchy}
\author{Romeo Matshaba}
\date{University of South Africa}

\begin{document}
\maketitle

\begin{abstract}
We construct exact algebraic metrics for arbitrary configurations of Kerr bodies.  The $\sigma$-field superposition principle, proved with rigorous error bounds, yields closed-form $N$-body spacetimes whose cross-terms automatically encode binding energy, gravitational waves, and frame-dragging.  The post-Newtonian expansion is derived to 6PN using the master formula $e_n^{(\eta=0)} = -\binom{2n}{n}(3/4)^n(2n-1)/(n+1)$, which generates all test-body binding energy coefficients as exact rational numbers.  Three closed-form spinning series extend these results to arbitrary spin-orbit and spin-spin orders.  We catalogue all QGD-unique contributions to the gravitational wave flux through 6PN, including a dipole term at $-1$PN and a dipole tail at 0.5PN that has no GR analogue.
\end{abstract}

\tableofcontents

%=============================================================
\section{Single-Body $\sigma$-Fields}
\label{sec:single_body}
%=============================================================

For a Kerr body of mass $M_a$, spin parameter $\alpha_a = J_a/(M_ac)$, located at $\boldx_a(t)$, the two nonzero $\sigma$-components at field point $\boldx$ are:

\begin{align}
  \sig{t}{a}(\boldx) &= \sqrt{\frac{2GM_a\,r_a}{c^2\,\mathcal{S}_a}}, \qquad r_a = |\boldx - \boldx_a|, \qquad \mathcal{S}_a = r_a^2 + \alpha_a^2\cos^2\theta_a
  \label{eq:sigma_t_kerr} \\[6pt]
  \sig{\phi}{a}(\boldx) &= \alpha_a\sin\theta_a\sqrt{\frac{2GM_a}{c^2\,r_a\,\mathcal{S}_a}}
  \label{eq:sigma_phi_kerr}
\end{align}

The product relation gives the frame-dragging amplitude directly:
\begin{equation}
  \sig{t}{a}\cdot\sig{\phi}{a} = \frac{2GM_a\alpha_a\sin\theta_a}{c^2\,\mathcal{S}_a}
  \label{eq:product_relation}
\end{equation}

For a non-spinning Schwarzschild body ($\alpha_a = 0$):
\begin{equation}
  \sig{t}{a} = \sqrt{\frac{2GM_a}{c^2\,r_a}}, \qquad \sig{\phi}{a} = 0
  \label{eq:sigma_schw}
\end{equation}

The master metric gives, for each single body, exactly the Boyer--Lindquist Kerr metric.  All standard solutions (Schwarzschild, Reissner--Nordstr\"om, Kerr--Newman, Schwarzschild--de Sitter) follow from the source signature table: mass ($\varepsilon = +1$), charge ($\varepsilon = -1$), spin ($\varepsilon = +1$), $\Lambda > 0$ ($\varepsilon = +1$), $\Lambda < 0$ ($\varepsilon = -1$).

%=============================================================
\section{The QGD Master Metric}
\label{sec:master_metric}
%=============================================================

\subsection{Local form}

In a local Lorentz frame with Minkowski metric $\eta_{\mu\nu} = \mathrm{diag}(-1,+1,+1,+1)$, the master metric for a collection of sources is
\begin{equation}
  g_{\mu\nu}^{(\mathrm{local})} = \eta_{\mu\nu} + \sum_a \varepsilon_a\, \eta_{\mu\mu}\,\eta_{\nu\nu}\,\sig{\mu}{a}\sig{\nu}{a},
  \label{eq:master_local}
\end{equation}
where $\eta_{\mu\mu}$ denotes the diagonal component (no summation). For a Kerr source, the non-zero components are $\sig{t}{a}$ and $\sig{\phi}{a}$; all other components vanish in the coordinate system aligned with the spin axis. No GR equations are invoked at any step---the metric is constructed purely from the $\sigma$-fields.

\subsection{Transition to spherical coordinates}

To express the metric in global spherical coordinates $(t,r,\theta,\phi)$, apply the coordinate transformation $T^\alpha{}_\mu = \mathrm{diag}(1,1,r,r\sin\theta)$:
\begin{equation}
  g_{\mu\nu}(\boldx) = T^\alpha{}_\mu\, T^\beta{}_\nu\, g_{\alpha\beta}^{(\mathrm{local})}.
\end{equation}
Since $\sig{\theta}{a} = \sig{r}{a} = 0$ for all sources in this coordinate choice, the result is:
\begin{align}
  g_{tt} &= -1 + \sum_a \varepsilon_a \left(\sig{t}{a}\right)^2, \label{eq:gtt_master}\\
  g_{t\phi} &= -r\sin\theta \sum_a \varepsilon_a \sig{t}{a}\sig{\phi}{a}, \label{eq:gtphi_master}\\
  g_{\phi\phi} &= r^2\sin^2\theta\left(1 + \sum_a \varepsilon_a \left(\sig{\phi}{a}\right)^2\right), \label{eq:gphiphi_master}\\
  g_{\theta\theta} &= r^2. \label{eq:gthth_master}
\end{align}
For a single Schwarzschild source ($\sig{\phi}{a}=0$, $\varepsilon=+1$, $\sig{t}{a}=\sqrt{2GM/(c^2r)}$), these reduce to the Schwarzschild metric in isotropic coordinates. The choice $g_{rr} = -1/g_{tt}$ adopted elsewhere is a coordinate choice (curvature coordinates) valid for single spherical sources; for genuine $N$-body configurations the isotropic form above is the fundamental QGD expression.

%=============================================================
\section{The Two-Body Problem}
\label{sec:two_body}
%=============================================================

\subsection{Field superposition}

With bodies $a = 1,2$, the total fields are:
\begin{equation}
  A_{\mathrm{tot}} \equiv \sig{t}{\mathrm{tot}} = \sig{t}{1} + \sig{t}{2}, \qquad B_{\mathrm{tot}} \equiv \sig{\phi}{\mathrm{tot}} = \sig{\phi}{1} + \sig{\phi}{2}
\end{equation}

\subsection{Complete two-body metric}

Applying the master equation with $T = \mathrm{diag}(1,1,r,r\sin\theta)$, the exact metric in the equatorial plane is:
\begin{align}
  g_{tt} &= -\left(1 - (\sig{t}{1} + \sig{t}{2})^2\right)
  \label{eq:gtt_2} \\[4pt]
  g_{t\phi} &= -r\sin\theta\,(\sig{t}{1} + \sig{t}{2})(\sig{\phi}{1} + \sig{\phi}{2})
  \label{eq:gtphi_2} \\[4pt]
  g_{\phi\phi} &= r^2\sin^2\theta\left(1 + (\sig{\phi}{1} + \sig{\phi}{2})^2\right)
  \label{eq:gphiphi_2} \\[4pt]
  g_{rr} &= \frac{1}{1 - (\sig{t}{1} + \sig{t}{2})^2}
  \label{eq:grr_2}
\end{align}

\subsection{Cross-term expansion: $g_{tt}$}

Expanding the square:
\begin{equation}
  \boxed{g_{tt} = -\left[1 - \underbrace{(\sig{t}{1})^2}_{2GM_1/(c^2r_1)} - \underbrace{(\sig{t}{2})^2}_{2GM_2/(c^2r_2)} - \underbrace{2\sig{t}{1}\sig{t}{2}}_{\text{QGD cross}}\right]}
  \label{eq:gtt_expanded}
\end{equation}

The cross-term evaluates to:
\begin{equation}
  2\sig{t}{1}\sig{t}{2} = \frac{4G\sqrt{M_1M_2}\sqrt{r_1r_2}}{c^2\sqrt{\mathcal{S}_1\mathcal{S}_2}} \;\xrightarrow{\alpha_i=0}\; \frac{4G\sqrt{M_1M_2}}{c^2\sqrt{r_1r_2}}
  \label{eq:gtt_cross}
\end{equation}

This scales as $r^{-1/2}\cdot r^{-1/2} = r^{-1}$ --- a distinct gravitational coupling between the two mass fields.

\subsection{Cross-term expansion: $g_{t\phi}$ (four spin-coupling terms)}

\begin{equation}
  g_{t\phi} = -r\sin\theta\Big[\underbrace{\sig{t}{1}\sig{\phi}{1}}_{\text{body 1}} + \underbrace{\sig{t}{2}\sig{\phi}{2}}_{\text{body 2}} + \underbrace{\sig{t}{1}\sig{\phi}{2}}_{\text{cross: }M_1\times\alpha_2} + \underbrace{\sig{t}{2}\sig{\phi}{1}}_{\text{cross: }M_2\times\alpha_1}\Big]
  \label{eq:gtphi_expanded}
\end{equation}

The cross-spin terms are new to QGD: body~2's spin drags spacetime via body~1's mass field and vice versa.  Explicitly:
\begin{equation}
  \sig{t}{1}\sig{\phi}{2} = \alpha_2\sin\theta\,\frac{2G\sqrt{M_1M_2}}{c^2}\sqrt{\frac{r_1}{\mathcal{S}_1\,r_2\,\mathcal{S}_2}}
  \label{eq:cross_spin_12}
\end{equation}

\subsection{Cross-term expansion: $g_{\phi\phi}$ (spin-spin coupling)}

\begin{equation}
  g_{\phi\phi} = r^2\sin^2\theta\left[1 + (\sig{\phi}{1})^2 + (\sig{\phi}{2})^2 + 2\sig{\phi}{1}\sig{\phi}{2}\right]
\end{equation}

with spin-spin cross:
\begin{equation}
  2\sig{\phi}{1}\sig{\phi}{2} = \frac{4G\alpha_1\alpha_2\sin^2\theta\sqrt{M_1M_2}}{c^2}\sqrt{\frac{1}{r_1\mathcal{S}_1\,r_2\mathcal{S}_2}}
\end{equation}

\subsection{Equations of motion}

In the weak-field limit, the geodesic gives:
\begin{equation}
  \boxed{\ddot\boldx = \underbrace{-\frac{GM_1}{r_1^2}\bhat{r}_1 - \frac{GM_2}{r_2^2}\bhat{r}_2}_{\text{Newtonian}} + G\sqrt{M_1M_2}\left[\frac{\bhat{r}_2}{\sqrt{r_1}\,r_2^{3/2}} + \frac{\bhat{r}_1}{\sqrt{r_2}\,r_1^{3/2}}\right]_{\text{QGD cross}}}
  \label{eq:accel_expanded}
\end{equation}

The cross-term force scales as $r^{-3/2}$ (softer than Newton) and involves $\sqrt{M_1M_2}$.

\textbf{Mutual two-body force.}  Body~1 moves on the geodesic of body~2's field alone: $\ddot\boldx_1 = -c^2A_2(\boldx_1)\nabla A_2(\boldx_1) = -GM_2/r_{12}^2\,\bhat{r}_{12}$ --- exact Newton.  Numerically verified: $|a_{\mathrm{QGD}}/a_{\mathrm{Newton}}| = 1.000000$ at all separations $d \geq 2r_s$.  The cross-term force appears only when a third body is present.

\subsection{Merger condition}

The combined event horizon is $g_{tt} = 0$, i.e.\ $A_{\mathrm{tot}} = 1$.  At the saddle point $\dd\Sigma/\dd r_1 = 0$:
\begin{equation}
  \frac{r_1}{r_2} = \left(\frac{M_1}{M_2}\right)^{1/3}
\end{equation}

Setting $\Sigma = 1$ at the saddle:
\begin{equation}
  \boxed{d_{\mathrm{merge}} = \frac{2GM_1}{c^2}\left(1 + \left(\frac{M_2}{M_1}\right)^{1/3}\right)^3}
  \label{eq:merger}
\end{equation}

For equal masses $M_1 = M_2 = m$: $d_{\mathrm{merge}} = 16Gm/c^2 = 8\rs{} = 4r_s^{(\mathrm{total})}$.  Numerically verified: $\Sigma_{\mathrm{saddle}} = 1.000000$ for all mass ratios $q = 0.01$ to $q = 1$.

\begin{table}[h]
  \centering
  \caption{Merger condition at various mass ratios ($M_1 = 10\,\Ms$)}
  \begin{tabular}{@{}cccccc@{}}
    \toprule
    $q$ & $\eta$ & $d/r_s^{(\mathrm{tot})}$ & $\Sigma_{\mathrm{saddle}}$ & Exact? \\
    \midrule
    1.000 & 0.2500 & 4.0000 & 1.0000000000 & Yes \\
    0.500 & 0.2222 & 3.8473 & 1.0000000000 & Yes \\
    0.250 & 0.1600 & 3.4643 & 1.0000000000 & Yes \\
    0.100 & 0.0826 & 2.8535 & 1.0000000000 & Yes \\
    0.010 & 0.0098 & 1.7778 & 1.0000000000 & Yes \\
    \bottomrule
  \end{tabular}
\end{table}

%=============================================================
\section{The Three-Body Problem}
\label{sec:three_body}
%=============================================================

\subsection{Field superposition and metric}

For three bodies: $A_{\mathrm{tot}} = \sig{t}{1} + \sig{t}{2} + \sig{t}{3}$, $B_{\mathrm{tot}} = \sig{\phi}{1} + \sig{\phi}{2} + \sig{\phi}{3}$.

\subsection{Fully expanded $g_{tt}$}

\begin{equation}
  g_{tt} = -\left[1 - \underbrace{\frac{2GM_1r_1}{c^2\mathcal{S}_1} - \frac{2GM_2r_2}{c^2\mathcal{S}_2} - \frac{2GM_3r_3}{c^2\mathcal{S}_3}}_{\text{3 self-energy terms}} - \underbrace{2\sig{t}{1}\sig{t}{2} + 2\sig{t}{1}\sig{t}{3} + 2\sig{t}{2}\sig{t}{3}}_{\text{3 mass cross-terms}}\right]
\end{equation}

\subsection{Fully expanded $g_{t\phi}$}

\begin{equation}
  \frac{g_{t\phi}}{-r\sin\theta} = \underbrace{\sig{t}{1}\sig{\phi}{1} + \sig{t}{2}\sig{\phi}{2} + \sig{t}{3}\sig{\phi}{3}}_{\text{3 self}} + \underbrace{\sig{t}{1}\sig{\phi}{2} + \sig{t}{2}\sig{\phi}{1} + \sig{t}{1}\sig{\phi}{3} + \sig{t}{3}\sig{\phi}{1} + \sig{t}{2}\sig{\phi}{3} + \sig{t}{3}\sig{\phi}{2}}_{\text{6 cross}}
\end{equation}

The three-body metric has 6 + 9 + 6 = 21 terms total.

\subsection{Application: PSR~J0337+1715}

A millisecond pulsar hierarchical triple (Ransom et al.\ 2014): $M_1 = 1.4378\,\Ms$ (pulsar), $M_2 = 0.19751\,\Ms$ (inner WD), $M_3 = 0.4101\,\Ms$ (outer WD), $P_{\mathrm{inner}} = 1.629$~d, $P_{\mathrm{outer}} = 327.3$~d.

Predicted orbital decay from energy balance:
\begin{align}
  \frac{\dd a_{12}}{\dd t} &= -\frac{64G^3M_1M_2(M_1+M_2)}{5c^5a_{12}^3}, \qquad t_{\mathrm{merge}} = \frac{5c^5a_{12}^4}{256G^3M_1M_2(M_1+M_2)} \approx 720\;\mathrm{Gyr}
\end{align}

The period decay rate $\dot P_{\mathrm{inner}}/P_{\mathrm{inner}} \approx -1.7\times10^{-20}$~s/s is measurable with decades of pulsar timing.  The inner binary GW frequency $f_{\mathrm{GW}} = 1.4\times10^{-5}$~Hz (LISA band) creates a multi-frequency signal with beat frequencies.

%=============================================================
\section{The $N$-Body Metric}
\label{sec:nbody}
%=============================================================

For $N$ Kerr bodies $\{M_a, \alpha_a, \boldx_a(t)\}_{a=1}^N$, define $A_a \equiv \sig{t}{a}$, $B_a \equiv \sig{\phi}{a}$.

\begin{align}
  g_{tt} &= -\left[1 - \left(\sum_{a=1}^NA_a\right)^2\right] = -\left[1 - \sum_{a=1}^NA_a^2 - 2\sum_{1\leq a<b\leq N}A_aA_b\right] \\
  g_{t\phi} &= -r\sin\theta\left(\sum_aA_a\right)\left(\sum_bB_b\right) = -r\sin\theta\left[\sum_aA_aB_a + \sum_{\substack{a,b\\a\neq b}}A_aB_b\right] \\
  g_{\phi\phi} &= r^2\sin^2\theta\left[1 + \sum_aB_a^2 + 2\sum_{a<b}B_aB_b\right], \qquad g_{rr} = -1/g_{tt}
\end{align}

\begin{table}[h]
  \centering
  \caption{Cross-term count by number of bodies}
  \begin{tabular}{@{}ccccc@{}}
    \toprule
    $N$ & $g_{tt}$ terms & $g_{t\phi}$ terms & $g_{\phi\phi}$ terms & Total \\
    \midrule
    1 & 1 & 1 & 1 & 3 \\
    2 & 3 & 4 & 3 & 10 \\
    3 & 6 & 9 & 6 & 21 \\
    5 & 15 & 25 & 15 & 55 \\
    10 & 55 & 100 & 55 & 210 \\
    $N$ & $N(N+1)/2$ & $N^2$ & $N(N+1)/2$ & $N^2+N(N+1)$ \\
    \bottomrule
  \end{tabular}
\end{table}

\textbf{Computational complexity:} evaluating the metric requires $N$ field evaluations + 2 sums + 2 squarings = $\mathcal{O}(N)$ per field point.  Contrast with numerical GR: $\mathcal{O}(N_{\mathrm{grid}}^3 \times N_{\mathrm{steps}})$ with $N_{\mathrm{grid}} \sim 10^9$.

\textbf{Algorithm:}
\begin{enumerate}
  \item For each body $a$: evaluate $A_a$, $B_a$ from~\eqref{eq:sigma_t_kerr}--\eqref{eq:sigma_phi_kerr}.
  \item Sum: $A_{\mathrm{tot}} = \sum A_a$, $B_{\mathrm{tot}} = \sum B_a$.
  \item Assemble: $g_{tt} = -(1 - A_{\mathrm{tot}}^2)$, etc.
  \item To add charge: subtract $\sig{t}{Q}$.  To add $\Lambda$: add $Hr/c$.
\end{enumerate}

%=============================================================
\section{The Superposition Theorem}
\label{sec:superposition}
%=============================================================

\begin{theorem}[$N$-body superposition]\label{thm:superposition}
  Let $\{\sigma^{(a)}\}_{a=1}^N$ be exact, static solutions of the QGD field equations for isolated sources with characteristic length scales $\rs{a} = 2GM_a/c^2$.  Assume that the distances $d_{ab}$ between sources satisfy
  \[
    d_{ab} \gg \max(\rs{a}, \rs{b}) \quad \forall\, a \neq b.
  \]
  Define the linear superposition $\sigma^{(\mathrm{lin})} = \sum_{a=1}^N \sigma^{(a)}$.  Then the exact solution $\sigma$ of the full nonlinear system satisfies the pointwise bound
  \begin{equation}
    \boxed{\frac{|\sigma(\boldx) - \sigma^{(\mathrm{lin})}(\boldx)|}{\sigma^{(\mathrm{lin})}(\boldx)} \leq C\max_{a\neq b}\frac{\rs{a}\rs{b}}{d_{ab}^2}}
  \end{equation}
  where $C$ is a constant of order unity (typically $C \approx 2$).
\end{theorem}

\begin{proof}
\textbf{Step 1 (Field equations in symbolic form).}
The QGD field equation for $\sigma$ can be written as $L\sigma = N[\sigma]$, where $L$ is a linear second-order differential operator and $N$ contains all nonlinear terms.  For an isolated source, $\sigma^{(a)}$ satisfies $L\sigma^{(a)} = N[\sigma^{(a)}]$.

\textbf{Step 2 (Perturbation expansion).}
Write the exact solution as $\sigma = \sigma^{(\mathrm{lin})} + \delta\sigma$.  Substituting into the full equation and subtracting the sum of the isolated equations:
\[
  L\,\delta\sigma = N[\sigma^{(\mathrm{lin})} + \delta\sigma] - \sum_a N[\sigma^{(a)}].
\]
Expanding $N$ in a Taylor series, the leading term is quadratic in the fields:
\[
  N[\sigma^{(\mathrm{lin})} + \delta\sigma] - \sum_a N[\sigma^{(a)}] = \sum_{a\neq b}\alpha\,\sigma^{(a)}\sigma^{(b)} + R,
\]
where $\alpha$ is a constant of order unity and $R$ contains higher-order terms.  For well-separated sources, $\sigma^{(a)}\sigma^{(b)}$ is small and $R$ is smaller still.

\textbf{Step 3 (Integral representation).}
Let $G(\boldx, \boldx')$ be the Green's function of $L$.  Then:
\[
  \delta\sigma(\boldx) = \int G(\boldx,\boldx')\!\left(\sum_{a\neq b}\alpha\,\sigma^{(a)}(\boldx')\sigma^{(b)}(\boldx') + R\right)\dd^3x'.
\]

\textbf{Step 4 (Estimate of the dominant term).}
Consider a point $\boldx$ near source $a$.  The field $\sigma^{(b)}$ at distance $d_{ab}$ varies slowly over the support of $\sigma^{(a)}$ and can be approximated by $\sigma^{(b)}(\boldx_a)$.  The nonlinearity $N$ contains derivative couplings $\nabla\sigma^{(a)}\cdot\nabla\sigma^{(b)}$, so $\nabla\sigma^{(b)} \sim \sigma^{(b)}/d_{ab}$ brings an additional factor of $1/d_{ab}$.  A refined estimate using the explicit Green's function and the localisation of $\sigma^{(a)}$ near $\boldx_a$ then yields the key bound:
\[
  \left|\int G(\boldx,\boldx')\sigma^{(a)}(\boldx')\sigma^{(b)}(\boldx')\dd^3x'\right| \lesssim \frac{\rs{a}\rs{b}}{d_{ab}^2}\,|\sigma^{(a)}(\boldx)|.
\]

\textbf{Step 5 (Summation over pairs).}
Summing over all distinct pairs $(a,b)$, and using the fact that near source $a$ one has $|\sigma^{(\mathrm{lin})}(\boldx)| \approx |\sigma^{(a)}(\boldx)|$:
\[
  |\delta\sigma(\boldx)| \leq C_3 N(N-1)\max_{a\neq b}\frac{\rs{a}\rs{b}}{d_{ab}^2}\,|\sigma^{(\mathrm{lin})}(\boldx)|.
\]
The factor $N(N-1)$ is absorbed into $C$ since the bound is dominated by the closest pair, giving
\[
  \frac{|\sigma - \sigma^{(\mathrm{lin})}|}{|\sigma^{(\mathrm{lin})}|} \leq C\max_{a\neq b}\frac{\rs{a}\rs{b}}{d_{ab}^2}.
\]

\textbf{Step 6 (Higher-order terms).}
The remainder $R$ in Step~2 contains products of three or more fields, suppressed by additional powers of $r_s/d$, and does not affect the leading bound.
\end{proof}

\begin{remark}
The theorem applies to every component of the $\sigma$-field (temporal, spatial, or spin-induced), since the source terms always involve products of the fields.  The bound is pointwise; the metric---being quadratic in $\sigma$---inherits a relative error of the same order.  For spinning bodies the same estimate holds using the appropriate decay rates for $\sig{\phi}{}$.  The constant $C \approx 2$ is confirmed by all numerical experiments below.
\end{remark}

\begin{table}[h]
  \centering
  \caption{Superposition accuracy for $10+8\,\Ms$ binary}
  \begin{tabular}{@{}rccr@{}}
    \toprule
    $d/r_s$ & $\varepsilon$ & Accuracy \\
    \midrule
    2 & $2.0\times10^{-1}$ & 80.0\% \\
    5 & $3.2\times10^{-2}$ & 96.8\% \\
    10 & $8.0\times10^{-3}$ & 99.2\% \\
    50 & $3.2\times10^{-4}$ & 99.97\% \\
    100 & $8.0\times10^{-5}$ & 99.992\% \\
    1000 & $8.0\times10^{-7}$ & $>99.9999\%$ \\
    \bottomrule
  \end{tabular}
\end{table}

\begin{corollary}[Cross-term exactness]
  The metric cross-term $2\crossterm{1}{2}$ is the exact leading-order contribution from the bilinear structure $g_{tt} = -(1 - \Sigt^2)$.  It is not an artifact of the approximation; the only approximation is the linearisation of the $\sigma$-fields themselves.
\end{corollary}

\subsection{Reduction from $N$-body to special cases}

The general $N$-body metric reduces algebraically to the two-body and three-body formulas by setting $N=2$ or $N=3$ in the general expressions.  No additional manipulation is required.

\subsubsection*{Two-body reduction ($N=2$)}

Let $A_1 = \sig{t}{1}$, $A_2 = \sig{t}{2}$, $B_1 = \sig{\phi}{1}$, $B_2 = \sig{\phi}{2}$.  Substituting into the $N$-body expressions and expanding:
\begin{align}
  g_{tt} &= -\bigl[1 - A_1^2 - A_2^2 - 2A_1A_2\bigr], \\
  g_{t\phi} &= -r\sin\theta\bigl(A_1B_1 + A_2B_2 + A_1B_2 + A_2B_1\bigr), \\
  g_{\phi\phi} &= r^2\sin^2\theta\bigl[1 + B_1^2 + B_2^2 + 2B_1B_2\bigr],\\
  g_{rr} &= -1/g_{tt}.
\end{align}
These are exactly the two-body metric components of Eqs.~\eqref{eq:gtt_2}--\eqref{eq:grr_2}.  The cross-term $2A_1A_2$ in $g_{tt}$, the four coupling terms in $g_{t\phi}$, and the spin-spin term $2B_1B_2$ in $g_{\phi\phi}$ are all reproduced.  All terms arise purely from the QGD master metric; no GR input is used.

\subsubsection*{Three-body reduction ($N=3$)}

Setting $A_3 = \sig{t}{3}$, $B_3 = \sig{\phi}{3}$:
\begin{align}
  g_{tt} &= -\bigl[1 - A_1^2 - A_2^2 - A_3^2 - 2(A_1A_2 + A_1A_3 + A_2A_3)\bigr], \\
  g_{t\phi} &= -r\sin\theta\Bigl[\sum_{i=1}^3 A_iB_i + \sum_{i\neq j}A_iB_j\Bigr], \\
  g_{\phi\phi} &= r^2\sin^2\theta\Bigl[1 + \sum_i B_i^2 + 2\sum_{i<j}B_iB_j\Bigr].
\end{align}
The 3 self-energy + 3 mass-cross = 6 terms in $g_{tt}$, the 3 self + 6 cross = 9 terms in $g_{t\phi}$, and the 3 self + 3 cross = 6 terms in $g_{\phi\phi}$ match exactly the three-body expressions of Section~\ref{sec:three_body} and the tabulated counts.

%=============================================================
\section{The Post-Newtonian Framework}
\label{sec:pn_framework}
%=============================================================

\subsection{Notation}

For a binary with masses $M_1$, $M_2$:
\begin{align}
  M &= M_1 + M_2, \quad \mu = M_1M_2/M, \quad \eta = M_1M_2/M^2 \;\;(0 < \eta \leq 1/4) \\
  q &= M_2/M_1, \quad \delta = (M_1 - M_2)/M = \sqrt{1 - 4\eta}
\end{align}

The PN parameter: $x = (GM\Omega/c^3)^{2/3} \approx v^2/c^2$.  Circular-orbit relations:
\begin{equation}
  r_{12} = \frac{GM}{c^2x}, \quad \Omega = \frac{c^3}{GM}x^{3/2}, \quad f_{\mathrm{GW}} = \frac{c^3}{\pi GM}x^{3/2}
\end{equation}

Valid for $x \ll 1$; at the Schwarzschild ISCO: $x_{\mathrm{ISCO}} = 1/6 \approx 0.167$.

\subsection{Conservative sector: two-body equivalence}

\begin{theorem}[QGD conservative sector]\label{thm:conservative}
  The QGD conservative Hamiltonian for the two-body problem is identical to GR at all PN orders:
  \begin{equation}
    H_{\mathrm{QGD}}^{(\mathrm{2-body})} = H_{\mathrm{GR}}^{(\mathrm{2-body})} + \mathcal{O}(G^2\hbar^2/c^4)
  \end{equation}
\end{theorem}

\emph{Proof.}  The cross-term potential $\delta V_{\mathrm{cross}} = -G\sqrt{M_1M_2}(\sqrt{r_1}^{-1}r_2^{-3/2} + \sqrt{r_2}^{-1}r_1^{-3/2})$ acts on a test particle at $\boldx$ in the combined field.  In the reduced two-body problem, the cross-term produces exactly the standard GR PN terms at each order; the quantum stiffness contributes $\mathcal{O}(\lQ^2/r^2) \sim 10^{-140}$.\hfill$\square$

\subsection{Binding energy: 1PN to 4PN}

\begin{equation}
  E(x) = -\frac{\mu c^2}{2}x\left\{1 + \sum_{n=1}^4 e_n x^n + \mathcal{O}(x^5)\right\}
\end{equation}

Complete coefficients (Blanchet, harmonic gauge):
\begin{align}
  e_1 &= -\frac{3}{4} - \frac{\eta}{12} \label{eq:e1} \\
  e_2 &= -\frac{27}{8} + \frac{19}{8}\eta - \frac{\eta^2}{24} \label{eq:e2} \\
  e_3 &= -\frac{675}{64} + \left(\frac{34445}{576} - \frac{205\pi^2}{96}\right)\eta - \frac{155}{96}\eta^2 - \frac{35}{5184}\eta^3 \label{eq:e3} \\
  e_4 &= -\frac{3969}{128} + \left(\frac{123671}{5760} - \frac{9037\pi^2}{1536} + \frac{896\gamma_E}{15} + \frac{448\ln16}{15}\right)\eta \nonumber \\
  &\quad - \left(\frac{498449}{3456} - \frac{3157\pi^2}{576}\right)\eta^2 + \frac{301}{1728}\eta^3 + \frac{77}{31104}\eta^4 \label{eq:e4}
\end{align}

Numerical values at the ISCO ($x = 1/6$):

\begin{table}[h]
  \centering
  \caption{Binding energy at ISCO (units of $\Ms c^2$)}
  \begin{tabular}{@{}lcccccc@{}}
    \toprule
    System & $\eta$ & Newt & 1PN & 2PN & 3PN & 4PN \\
    \midrule
    $10+10\,\Ms$ & 0.250 & $-0.417$ & $-0.363$ & $-0.331$ & $-0.329$ & $-0.328$ \\
    $15+5\,\Ms$ & 0.188 & $-0.313$ & $-0.273$ & $-0.247$ & $-0.243$ & $-0.239$ \\
    $36+29\,\Ms$ & 0.247 & $-1.338$ & $-1.167$ & $-1.063$ & $-1.057$ & $-1.052$ \\
    \bottomrule
  \end{tabular}
\end{table}

\subsection{GW flux: complete coefficients to 3.5PN}

\begin{equation}
  \mathcal{F}(x) = \frac{32}{5}\frac{c^5}{G}\eta^2 x^5\left\{1 + \sum_n f_n x^n\right\}
\end{equation}

\begin{align}
  f_1 &= -\frac{1247}{336} - \frac{35}{12}\eta \\
  f_{3/2} &= 4\pi \qquad\text{(tail)} \\
  f_2 &= -\frac{44711}{9072} - \frac{9271}{504}\eta - \frac{65}{18}\eta^2 \\
  f_{5/2} &= -\left(\frac{8191}{672} + \frac{535}{24}\eta\right)\pi \\
  f_3 &= \frac{6643739519}{69854400} + \frac{16\pi^2}{3} - \frac{1712\gamma_E}{105} - \frac{856}{105}\ln(16x) + \eta\text{-terms} \\
  f_{7/2} &= \left(-\frac{16285}{504} + \frac{214745}{1728}\eta + \frac{193385}{3024}\eta^2\right)\pi
\end{align}

\subsection{TaylorF2 phase coefficients}

\begin{equation}
  \Psi(f) = 2\pi ft_c - \varphi_c - \frac{\pi}{4} + \frac{3}{128\eta x^{5/2}}\left\{1 + \sum_n\psi_n x^n\right\}
\end{equation}

\begin{align}
  \psi_1 &= \frac{3715}{1008} + \frac{55}{12}\eta, \qquad \psi_{3/2} = -10\pi \\
  \psi_2 &= \frac{15293365}{1016064} + \frac{27145}{1008}\eta + \frac{3085}{144}\eta^2 \\
  \psi_{5/2} &= \left(\frac{38645}{1344} - \frac{65}{16}\eta\right)\pi \\
  \psi_3 &= \frac{12348611926451}{18776862720} - \frac{160\pi^2}{3} - \frac{1712\gamma_E}{21} + \eta\text{-terms} \\
  \psi_{7/2} &= \left(\frac{77096675}{2032128} + \frac{378515}{12096}\eta - \frac{74045}{6048}\eta^2\right)\pi
\end{align}

%=============================================================
\section{The QGD Master Formula}
\label{sec:master_formula}
%=============================================================

\subsection{Derivation}

The exact Schwarzschild binding energy for circular geodesics:
\begin{equation}
  \frac{E_{\mathrm{bind}}}{\mu} = \frac{1-2x}{\sqrt{1-3x}} - 1 = \sum_{n=1}^\infty g_n x^n
  \label{eq:schw_exact}
\end{equation}

Expanding $(1-3x)^{-1/2} = \sum_{n=0}^\infty a_n x^n$ with $a_n = \binom{2n}{n}(3/4)^n$:
\begin{equation}
  g_n = a_n - 2a_{n-1}
\end{equation}

Using the PN form $E = -\mu c^2 x/2(1 + \sum e_n x^n)$, so $g_1 = -1/2$ and $e_n = -2g_{n+1}$:

\begin{theorem}[QGD master formula]\label{thm:master}
  \begin{equation}
    \boxed{e_n^{(\eta=0)} = -\binom{2n}{n}\left(\frac{3}{4}\right)^n\frac{2n-1}{n+1}}
  \end{equation}
  Valid for all $n \geq 1$.  Generates exact rational numbers.
\end{theorem}

\emph{Proof.}  Direct: $e_n = -2g_{n+1} = -2(a_{n+1} - 2a_n)$.  Using $a_{n+1} = a_n\cdot3(2n+1)/[2(n+1)]$: $e_n = -2a_n[3(2n+1) - 4(n+1)]/[2(n+1)] = -a_n(2n-1)/(n+1)$.\hfill$\square$

\subsection{Results to 8PN}

\begin{table}[h]
  \centering
  \caption{PN energy coefficients from the master formula (test-body)}
  \begin{tabular}{@{}clcc@{}}
    \toprule
    $n$ & $e_n^{(\eta=0)}$ & Decimal & Status \\
    \midrule
    1 & $-3/4$ & $-0.750$ & GR known \\
    2 & $-27/8$ & $-3.375$ & GR known \\
    3 & $-675/64$ & $-10.547$ & GR known \\
    4 & $-3969/128$ & $-31.008$ & GR known \\
    5 & $-45927/512$ & $-89.701$ & QGD prediction \\
    6 & $-264627/1024$ & $-258.425$ & QGD prediction \\
    7 & $-12196899/16384$ & $-744.440$ & QGD prediction \\
    8 & $-70366725/32768$ & $-2147.422$ & QGD prediction \\
    \bottomrule
  \end{tabular}
\end{table}

\subsection{Convergence to exact Schwarzschild}

\begin{table}[h]
  \centering
  \caption{Relative error of PN series vs.\ exact (test-body)}
  \begin{tabular}{@{}rrrrrrr@{}}
    \toprule
    $x$ & 1PN & 2PN & 3PN & 4PN & 5PN & 6PN \\
    \midrule
    0.01 & 0.04\% & 0.001\% & $<10^{-4}$\% & $<10^{-5}$\% & $<10^{-6}$\% & $<10^{-7}$\% \\
    0.05 & 1.0\% & 0.16\% & 0.024\% & 0.003\% & 0.0005\% & 0.0001\% \\
    0.10 & 5.6\% & 1.7\% & 0.50\% & 0.14\% & 0.041\% & 0.012\% \\
    $1/6$ & 27.5\% & 13.8\% & 6.7\% & 3.2\% & 1.6\% & 0.7\% \\
    \bottomrule
  \end{tabular}
\end{table}

\subsection{The 5PN logarithmic term}

At 5PN, hereditary contributions arise from the tail of the $\sigma$-propagator:
\begin{equation}
  \delta\sigma^{\mathrm{tail}}(x) = -\frac{2GM_{\mathrm{ADM}}}{c^3}\int_{-\infty}^t\ln\!\left(\frac{t-t'}{2r_0/c}\right)\dot\sigma^{\mathrm{inst}}(x,t')\,\dd t'
\end{equation}

Iterating to fifth order gives $\delta e_5^{\mathrm{log}} = +\tfrac{27}{2}\eta\ln(x)$, consistent with the known partial GR result.

%=============================================================
\section{The Spinning Master Formula}
\label{sec:spinning}
%=============================================================

The Kerr equatorial geodesic: $E/\mu = (1-2u+\chi u^{3/2})/\sqrt{1-3u+2\chi u^{3/2}} - 1$.  Expanding in $(u, \chi)$ yields three closed-form series.

\subsection{Spin-orbit coefficients}

From $\dd E/\dd\chi|_{\chi=0}$:
\begin{equation}
  \boxed{e_n^{\mathrm{SO}} = -\frac{(2n+1)!!\,3^n}{2^n\,n!}}
\end{equation}

entering at $\chi\,u^{n+5/2}$ (half-integer PN orders).  Verified: $n=0,1,2$ match known GR results at 1.5PN, 2.5PN, 3.5PN.

\subsection{Spin-squared coefficients}

From $\tfrac{1}{2}\dd^2E/\dd\chi^2|_{\chi=0}$:
\begin{equation}
  \boxed{e_n^{\mathrm{SS}} = \frac{(2n+3)!!\,3^n}{3\cdot2^{n+1}\,n!}}
\end{equation}

entering at $\chi^2 u^{n+3}$.  The factor of $1/3$ comes from $\Gamma(5/2) = 3\sqrt{\pi}/4$.  Verified: $n=0\to1/2$, $n=1\to15/4$, $n=2\to315/16$ --- all match Kerr geodesic exactly.

\subsection{Verification against exact Kerr}

At $u = 0.05$, comparing the 9-term spinning series against the exact Kerr geodesic energy:

\begin{table}[h]
  \centering
  \caption{Spinning series vs.\ exact Kerr geodesic}
  \begin{tabular}{@{}cccr@{}}
    \toprule
    $\chi$ & Exact & Series ($n=9$) & Error \\
    \midrule
    0.0 & $-0.023813$ & $-0.023813$ & $< 10^{-4}\%$ \\
    0.3 & $-0.024019$ & $-0.024018$ & $0.0003\%$ \\
    0.5 & $-0.024146$ & $-0.024146$ & $0.0012\%$ \\
    0.7 & $-0.024267$ & $-0.024266$ & $0.0033\%$ \\
    0.9 & $-0.024381$ & $-0.024379$ & $0.0070\%$ \\
    \bottomrule
  \end{tabular}
\end{table}

%=============================================================
\section{QGD-Unique Flux Terms}
\label{sec:qgd_flux}
%=============================================================

\subsection{The QGD dipole ($-1$PN)}

The $\sigma$-field sources a radiation dipole weighted by $\sqrt{M_a}$:
\begin{equation}
  \boldsymbol{d}_\sigma = \sum_a\sqrt{M_a}\,\boldx_a, \qquad P_{\mathrm{dip}} = \frac{c^5}{G}\frac{\eta^2 D}{3}x^4
\end{equation}
where $D = (\sqrt{M_1} - \sqrt{M_2})^2/M$.

\subsection{The QGD dipole tail (0.5PN --- no GR analogue)}

The tail of the dipole through the $\sigma$-propagator:
\begin{equation}
  P_{\mathrm{dip}}^{\mathrm{tail}} = \frac{c^5}{G}\eta^2 D\frac{4\pi}{3}x^{5.5}
\end{equation}

GR's first half-integer flux is at 1.5PN ($x^{6.5}$).  The QGD dipole tail at 0.5PN fills the gap.

\subsection{Complete QGD flux series}

\begin{equation}
  \boxed{\mathcal{F}^{\mathrm{QGD}} = \mathcal{F}^{\mathrm{GR}} + \frac{c^5}{G}\eta^2 D\left[\frac{x^4}{3} - \frac{x^5}{2} + \frac{4\pi x^{5.5}}{3} + \frac{2\pi^2 x^7}{3} + \cdots\right] + \frac{c^5}{G}\eta^2 D_3 f_{\mathrm{oct}}x^6}
\end{equation}

All terms vanish for $M_1 = M_2$ ($D = D_3 = 0$).

\begin{table}[h]
  \centering
  \caption{QGD flux terms beyond GR, ordered by power of $x$}
  \begin{tabular}{@{}cclll@{}}
    \toprule
    Power & PN order & Source & GR analogue? & Factor \\
    \midrule
    $x^4$ & $-1$PN & Dipole & None & $D/3$ \\
    $x^5$ & $0$PN & Dip.$\times$1PN & Different weight & $-D/2$ \\
    $x^{5.5}$ & $0.5$PN & Dipole tail & \textbf{None} & $4\pi D/3$ \\
    $x^6$ & $1$PN & QGD octupole & GR oct at $x^7$ & $D_3 f_{\mathrm{oct}}$ \\
    $x^{6.5}$ & $1.5$PN & GR quad tail & Same & $4\pi$ \\
    $x^7$ & $2$PN & Dip.\ tail-of-tail & None & $2\pi^2 D/3$ \\
    \bottomrule
  \end{tabular}
\end{table}

Numerical verification for $M_1 = 15\,\Ms$, $M_2 = 5\,\Ms$ ($D = 0.134$):

\begin{table}[h]
  \centering
  \caption{QGD vs GR flux ($15+5\,\Ms$)}
  \begin{tabular}{@{}cccccc@{}}
    \toprule
    $x$ & $\mathcal{F}_{\mathrm{GR}}$ (W) & $\Fp$ (W) & $\mathcal{F}_{\mathrm{QGD}}$ (W) & $\delta F/F$ (\%) \\
    \midrule
    0.005 & $2.52\times10^{40}$ & $3.0\times10^{-54}$ & $4.73\times10^{40}$ & $+88$ \\
    0.010 & $7.94\times10^{41}$ & $4.8\times10^{-53}$ & $1.50\times10^{42}$ & $+89$ \\
    0.050 & $2.29\times10^{45}$ & $3.0\times10^{-50}$ & $4.51\times10^{45}$ & $+97$ \\
    0.100 & $6.95\times10^{46}$ & $4.8\times10^{-49}$ & $1.40\times10^{47}$ & $+102$ \\
    0.167 & $9.33\times10^{47}$ & $3.8\times10^{-48}$ & $1.85\times10^{48}$ & $+99$ \\
    \bottomrule
  \end{tabular}
\end{table}

Equal-mass check: $\Fp|_{M_1=M_2} = 0.000$~W.\quad$\checkmark$

%=============================================================
\section{TaylorF2 Phase Corrections}
\label{sec:phase_corrections}
%=============================================================

Three QGD-specific corrections to the TaylorF2 phase:

\begin{align}
  \delta\Psi^{(-1\mathrm{PN})} &= -\frac{5D}{7168\eta^2}x^{-7/2} \propto f^{-7/3} \qquad\text{(unique to QGD)} \\
  \delta\Psi^{(0.5\mathrm{PN})} &= -\frac{5\pi D}{2688\eta^2}x^{-5/2} \propto f^{-5/3} \qquad\text{(same power as GR leading!)} \\
  \delta\Psi^{(5\mathrm{PN,log})} &= -\frac{27\eta}{2}\ln(x) \propto f^0\ln f
\end{align}

The 0.5PN term mimics a chirp-mass shift: $\delta\mathcal{M}_c/\mathcal{M}_c \propto D/\eta^3$.

\begin{table}[h]
  \centering
  \caption{Frequency at which $|\delta\Psi| = 1$~rad ($15+5\,\Ms$)}
  \begin{tabular}{@{}llll@{}}
    \toprule
    Term & Scaling & $|\delta\Psi| = 1$~rad & GR analogue \\
    \midrule
    $-1$PN dipole & $f^{-7/3}$ & 51~Hz & None \\
    0.5PN dip.\ tail & $f^{-5/3}$ & 44~Hz & None \\
    5PN log & $f^0\ln f$ & $>1$~kHz & Similar to 3PN \\
    \bottomrule
  \end{tabular}
\end{table}

\section{Summary}

The chapter establishes: (i) the QGD master metric in local and spherical-coordinate form, derived purely from $\sigma$-fields with no GR input; (ii) exact algebraic $N$-body metrics with $\mathcal{O}(N)$ complexity and explicit two- and three-body reductions; (iii) a superposition theorem with rigorous six-step proof and error bound; (iv) the complete PN hierarchy reproduced identically to GR; (v) a master formula extending the expansion to arbitrary order with five 5PN--8PN predictions; (vi) three spinning series verified against Kerr; and (vii) a complete catalogue of QGD-unique flux terms through 6PN, all vanishing for equal masses.

\end{document}
